\chapter{복합계, 얽힘, 그리고 싱글렛}

\lectureoverview{광자 방출에 관한 질문은 양자 사건과 앙상블 평균의 차이를 드러낸다. 이 구분에서 측정과 붕괴, 그리고 계를 측정 장치와 함께 다뤄야 할 필요성이 자연스럽게 이어진다. 따라서 두 계의 텐서곱 상태 공간을 구성하고, 이를 두 스핀의 경우로 구체화한다. 곱상태와 얽힌 상태를 구분하고, 각 한 스핀 관측량을 복합 공간에 포함한다. 이어서 싱글렛 상태를 결정적인 예로 살펴본다. 각 스핀은 국소적으로 편극되지 않았지만, 같은 방향 성분들은 완벽한 반상관을 보인다. 기저에 무관한 싱글렛의 형태를 살펴보면 벨 정리의 토대가 마련되며, 비고전적 상관관계와 광속보다 빠른 신호 전달의 차이도 분명해진다.}

\section{중첩된 스핀의 불연속적 전이}

먼저 \(z\)축 방향 자기장 속에 놓인 스핀 하나를 생각하자. \(\hbar=1\) 단위계에서
\begin{equation}
H=\frac{\omega}{2}\sigma_z,
\qquad
E_+=+\frac{\omega}{2},
\qquad
E_-=-\frac{\omega}{2}.
\label{eq:l6-spin-hamiltonian}
\end{equation}
두 에너지 준위의 간격은 다음과 같다.
\begin{equation}
\Delta E=E_+-E_-=\omega.
\label{eq:l6-energy-gap}
\end{equation}
복사장을 포함하면, 높은 준위에 있는 스핀은 낮은 준위로 내려오면서 에너지 \(\omega\)인 광자를 방출할 수 있다. 이미 낮은 준위에 있는 스핀에는 복사 전이를 일으킬 더 낮은 스핀 상태가 없다.

\begin{figure}[htbp]
\centering
\includegraphics[width=0.86\textwidth]{lecture_06_figure_01.png}
\caption{두 스핀 에너지 준위와 그 사이의 간격 \(\omega\), 여기서 \(\hbar=1\).}
\lectureframe{00:02:35}
\end{figure}

미묘한 경우는 옆 방향으로 준비된 스핀이다. 고전적인 자기 쌍극자 회전자는 자기장과 이루는 각의 코사인에 비례하는 에너지를 연속적으로 갖는다. 각주파수 \(\omega\)로 세차하면서 같은 진동수의 복사를 방출하고, 점차 에너지가 가장 낮은 방향으로 향한다. 방출되는 총 에너지는 초기 각도에 따라 연속적으로 달라진다. 처음에 수평인 회전자는 꼭대기에서 출발한 회전자 에너지의 절반을 방출한다.

양자역학에서 이 고전적 결과는 평균으로서만 성립한다. 다음과 같이 쓰자.
\begin{equation}
\ket{\psi}=\alpha\ket{u}+\beta\ket{d},
\qquad
|\alpha|^2+|\beta|^2=1,
\label{eq:l6-superposed-spin}
\end{equation}
여기서 \(\ket{u}\)와 \(\ket{d}\)는 각각 에너지가 높은 상태와 낮은 상태의 고유상태다. 높은 에너지 성분은 붕괴할 수 있지만 낮은 에너지 성분은 붕괴할 수 없다. 따라서
\begin{equation}
P(\text{one photon})=|\alpha|^2,
\qquad
P(\text{no photon})=|\beta|^2.
\end{equation}
광자가 방출된다면 그 에너지는 준위 간격 전체인 \(\omega\)이지, 그 일부가 아니다. 그러나 방출 에너지의 평균은 다음과 같다.
\begin{equation}
\expect{E_{\rm emitted}}=|\alpha|^2\omega.
\end{equation}
\(x\) 방향의 스핀에서는 \(|\alpha|^2=|\beta|^2=1/2\)이므로 평균은 \(\omega/2\)이다. 개별 사건에서는 여전히 에너지가 \(\omega\)인 광자 하나가 방출되거나 광자가 전혀 방출되지 않는다. 연속적인 고전적 결과는 앙상블 평균에 해당하고, 불연속성은 각각의 사건에 나타난다.

\begin{classroomqa}
\audiencequestion{흡수된 광자의 에너지는 준위 간격과 얼마나 정확히 일치해야 하나요?}
\lecturerresponse{공명선에는 유한한 선폭이 있습니다. 들뜬 상태의 수명이 \(\tau\)라면, 그 특성 에너지 폭은 대략 \(\Delta E\sim\hbar/\tau\)입니다. 일반적으로 자기 쌍극자 모멘트가 클수록 복사장과 더 강하게 결합하고, 수명은 짧아지며, 선폭은 넓어집니다. 이것이 불안정한 준위에 적용되는 에너지-시간 불확정성 관계입니다.}
\end{classroomqa}

\begin{classroomqa}
\audiencequestion{광자 방출 자체를 스핀을 측정하는 장치로 쓸 수 있나요?}
\lecturerresponse{붕괴가 일어났는지는 알 수 있지만, 여기서 필요한 이상적인 반복 측정 모형은 아닙니다. 스핀 성분을 측정했다면 기록된 값에 해당하는 고유상태에 스핀이 남아 있어야 합니다. 반면 붕괴가 일어난 각 가지는 모두 낮은 준위로 이행하므로, 가능한 두 입력값을 각각의 대응 고유상태를 보존한 채 단순히 기록하는 과정이 아닙니다.}
\end{classroomqa}

\section{붕괴와 복합계가 필요한 이유}

외부 개입이 없는 동안 고립계의 상태는 Schr\"odinger 방정식에 따라 진화한다. 측정은 보통 여기에 추가되는 규칙으로 제시된다. \(L\)의 정규직교 고유벡터가 \(\ket{\lambda_i}\)라고 하고, 상태를 다음과 같이 전개하자.
\begin{equation}
\ket{\psi}=\sum_i\alpha_i\ket{\lambda_i},
\qquad
L\ket{\lambda_i}=\lambda_i\ket{\lambda_i}.
\label{eq:l6-measurement-expansion}
\end{equation}
측정값 하나인 \(\lambda_k\)가 확률 \(|\alpha_k|^2\)로 나오며, 그 뒤 상태는 다음과 같다.
\begin{equation}
\ket{\psi}\longrightarrow\ket{\lambda_k}.
\label{eq:l6-collapse}
\end{equation}
측정 후의 기술에서 나머지 성분들이 사라지는 것을 붕괴라고 한다.

\begin{figure}[htbp]
\centering
\includegraphics[width=0.84\textwidth]{lecture_06_figure_02.png}
\caption{관측량 \(L\)의 고유기저로 전개한 상태. 하나의 결과가 선택되기 전이다.}
\lectureframe{00:12:35}
\end{figure}

\begin{classroomqa}
\audiencequestion{붕괴 규칙에 문제가 있나요?}
\lecturerresponse{측정 결과를 예측하는 조작적 규칙으로서는 문제가 없습니다. 개념적인 긴장은 측정 장치도 물리계라는 데서 생깁니다. 대상과 장치가 모두 양자역학을 따른다면, 두 계의 상호작용에 유니터리 Schr\"odinger 진화 외에 왜 두 번째 운동 법칙이 필요할까요? 여기서는 대상과 장치를 하나의 더 큰 양자계로 기술하는 방법을 택합니다.}
\end{classroomqa}

계와 측정 장치 사이에 명확한 경계를 긋는 문제는 역사적인 논쟁의 핵심이었다. 선언만으로 그 논쟁을 끝내기보다, 먼저 양자계를 어떻게 결합하는지 살펴보자. 결합된 상태 공간을 마련하면, 계와 장치 사이에 상관관계를 만드는 상호작용으로 측정을 연구할 수 있다. 따라서 당장 답해야 할 수학적 질문은 다음과 같다.
\begin{quote}
양자계 두 개가 주어졌을 때, 전체 계의 상태 공간은 무엇인가?
\end{quote}

\section{텐서곱 상태 공간}

두 계를 \(A\)와 \(I\)라 하고, 각각의 힐베르트 공간을 \(\mathcal S_A\)와 \(\mathcal S_I\)라고 하자. 복합계의 공간은 다음과 같다.
\begin{equation}
\mathcal S_{AI}=\mathcal S_A\otimes\mathcal S_I.
\label{eq:l6-tensor-space}
\end{equation}
두 인자 공간의 정규직교기저를 각각 \(\{\ket{a}\}\)와 \(\{\ket{i}\}\)로 택하자. 첫 번째 계의 각 기저벡터를 두 번째 계의 모든 기저벡터와 짝짓는다.
\begin{equation}
\ket{a}\otimes\ket{i}\equiv\ket{a\,i}.
\end{equation}
인자 공간의 차원이 각각 \(D_A\)와 \(D_I\)이면,
\begin{equation}
\dim\mathcal S_{AI}=D_A D_I.
\end{equation}
내적은 인자별로 분해된다.
\begin{equation}
\braket{b\,j}{a\,i}
=\braket{b}{a}\braket{j}{i}
=\delta_{ba}\delta_{ji}.
\label{eq:l6-tensor-inner-product}
\end{equation}
따라서 두 복합계 기저 상태는 어느 한 부분계의 표지가 다르기만 해도 서로 직교한다.

\begin{figure}[htbp]
\centering
\includegraphics[width=0.86\textwidth]{lecture_06_figure_03.png}
\caption{각 상태 공간과 그 텐서곱, 그리고 짝지은 정규직교기저.}
\lectureframe{00:24:05}
\end{figure}

짝지은 기저벡터만으로 공간 전체가 이루어지는 것은 아니다. 이 벡터들의 완전한 선형생성공간은 다음과 같다.
\begin{equation}
\ket{\Psi}=\sum_{a,i}c_{ai}\ket{a\,i},
\qquad
\sum_{a,i}|c_{ai}|^2=1.
\label{eq:l6-general-tensor-state}
\end{equation}

\begin{classroomqa}
\audiencequestion{\(\ket{a}\)와 \(\ket{i}\)는 임의의 벡터인가요?}
\lecturerresponse{여기서는 선택한 정규직교기저의 원소입니다. 각 인자 공간에서 어떤 정규직교기저를 택해도 됩니다. 그 기저벡터들을 짝지으면 복합계 공간의 정규직교기저가 만들어지고, 그 뒤에는 식~\eqref{eq:l6-general-tensor-state}처럼 선형결합을 취해 임의의 벡터를 얻습니다.}
\end{classroomqa}

\begin{classroomqa}
\audiencequestion{텐서곱은 선택한 기저에 따라 달라지나요?}
\lecturerresponse{좌표로 나타내는 방식은 달라지지만 공간 자체는 달라지지 않습니다. 어느 부분계에서 기저를 바꾸더라도 \(\mathcal S_A\otimes\mathcal S_I\)의 동일한 벡터를 나타내는 유도 기저만 바뀝니다.}
\end{classroomqa}

\section{두 스핀과 곱상태}

이제 스핀 두 개를 생각하자. 첫 번째 스핀의 성분은 \(\sigma_x,\sigma_y,\sigma_z\), 두 번째 스핀의 성분은 \(\tau_x,\tau_y,\tau_z\)로 나타내자. \(z\) 기저에서 복합계 기저는 다음과 같다.
\begin{equation}
\ket{uu},\qquad \ket{ud},\qquad \ket{du},\qquad \ket{dd}.
\label{eq:l6-two-spin-basis}
\end{equation}
첫 번째 자리는 항상 \(\sigma\), 두 번째 자리는 \(\tau\)를 가리킨다. 일반적인 상태는 다음과 같다.
\begin{equation}
\ket{\Psi}
=c_{uu}\ket{uu}+c_{ud}\ket{ud}
+c_{du}\ket{du}+c_{dd}\ket{dd}.
\label{eq:l6-general-two-spin}
\end{equation}

\begin{classroomqa}
\audiencequestion{두 스핀이 상호작용하면 기저벡터 자체도 바뀌나요?}
\lecturerresponse{동역학을 기술할 때 고정된 기저를 사용할 필요는 없습니다. 선택한 기저에서는 상호작용에 따라 계수 \(c_{ab}(t)\)가 바뀌지만, \(\ket{uu},\ket{ud},\ket{du},\ket{dd}\)의 정의는 바뀌지 않습니다. \(\ket{uu}\)로 시작한 상태는 진화하면서 중첩 상태, 심지어 \(\ket{dd}\)가 될 수도 있지만, 좌표 기저는 그대로 유지됩니다.}
\end{classroomqa}

이 구성은 일관성 검사도 통과해야 한다. 두 스핀이 멀리 떨어져 있고 \(\sigma\)만 살펴본다면, 예측은 \(\tau\)를 도입하기 전에 전개한 한 스핀 이론과 일치해야 한다. 텐서곱은 이 국소 이론을 보존하도록 고안되었다.

두 스핀 상태 중에는 서로 독립적으로 준비된 스핀을 나타내는 상태가 있다. 예를 들어 다음 상태를 생각하자.
\begin{equation}
\ket{R}=\frac{\ket{u}+\ket{d}}{\sqrt2}
\end{equation}
이는 \(+1\)인 \(x\) 성분의 고유상태다. 그러면
\begin{equation}
\ket{u}_{\sigma}\otimes\ket{R}_{\tau}
=\frac{\ket{uu}+\ket{ud}}{\sqrt2}.
\label{eq:l6-up-right-product}
\end{equation}

\begin{classroomqa}
\audiencequestion{두 인자를 서로 다른 축 방향으로 준비해도 되나요?}
\lecturerresponse{그렇습니다. 각 인자는 임의의 한 스핀 상태일 수 있습니다. 식~\eqref{eq:l6-up-right-product}에서 첫 번째 스핀은 \(z\) 방향으로, 두 번째 스핀은 \(x\) 방향으로 확정되어 있습니다. 두 번째 인자를 \(z\) 기저로 전개하는 것은 복합계 기저에서 같은 곱상태를 다시 쓰는 것일 뿐입니다.}
\end{classroomqa}

가장 일반적인 곱상태는 다음과 같다.
\begin{align}
\ket{\Psi_{\rm prod}}
&=(\alpha_u\ket{u}+\alpha_d\ket{d})
\otimes(\beta_u\ket{u}+\beta_d\ket{d})\\
&=\alpha_u\beta_u\ket{uu}
+\alpha_u\beta_d\ket{ud}
\\
&\quad
+\alpha_d\beta_u\ket{du}
+\alpha_d\beta_d\ket{dd}.
\label{eq:l6-product-expansion}
\end{align}

\begin{figure}[htbp]
\centering
\includegraphics[width=0.86\textwidth]{lecture_06_figure_04.png}
\caption{네 상태로 이루어진 복합계 기저로 전개한 인수분해된 두 스핀 상태.}
\lectureframe{00:37:45}
\end{figure}

\begin{classroomqa}
\audiencequestion{식~\eqref{eq:l6-product-expansion}에는 어떤 종류의 벡터곱이 나오나요?}
\lecturerresponse{전체 공간은 두 힐베르트 공간의 텐서곱입니다. 각 공간에서 가져온 벡터 하나씩의 곱으로 인수분해되는 벡터를 곱상태 또는 단순 텐서라고 합니다. 텐서곱 공간의 대부분의 벡터는 이런 곱들의 합이며, 하나의 곱으로 쓸 수 없습니다.}
\end{classroomqa}

계수에서 이 제약을 직접 확인할 수 있다.
\begin{equation}
c_{uu}c_{dd}=c_{ud}c_{du}
\qquad\text{for every product state}.
\label{eq:l6-rank-one-condition}
\end{equation}
다시 말해, 계수 행렬 \(2\times2\)의 계수는 1이다. 매개변수의 개수를 세어도 같은 결론에 이른다. 복소 계수 네 개에는 실수 매개변수가 여덟 개 있다. 정규화 조건으로 하나가, 무관한 전체 위상으로 하나가 제거되어 일반적인 순수 두 스핀 상태에는 여섯 개가 남는다. 순수 한 스핀 상태에는 실수 매개변수가 두 개 필요하므로, 별도로 정해지는 스핀 두 개에는 네 개만 필요하다. 따라서 곱상태는 여섯 매개변수로 이루어진 두 스핀 상태 공간 안의 특별한 네 매개변수 부분집합이다. 나머지 상태는 얽혀 있다.

\section{복합 공간의 국소 관측량}

익숙한 한 스핀 관측량은 더 큰 계 안에서도 그대로 성립해야 한다. 구체적으로,
\begin{equation}
\sigma_m\equiv\sigma_m\otimes I,
\qquad
\tau_n\equiv I\otimes\sigma_n,
\qquad
m,n\in\{x,y,z\}.
\label{eq:l6-local-operators}
\end{equation}
자리를 나타내는 약식 표기법을 정해 두면 간결한 표기를 써도 문제가 없다. 다음을 사용하면
\begin{equation}
\begin{aligned}
\sigma_x&=
\begin{pmatrix}0&1\\1&0\end{pmatrix},\\[0.4em]
\sigma_y&=
\begin{pmatrix}0&-i\\i&0\end{pmatrix},\\[0.4em]
\sigma_z&=
\begin{pmatrix}1&0\\0&-1\end{pmatrix},
\end{aligned}
\end{equation}
첫 번째 스핀에 작용하는 연산자는 두 번째 스핀의 표지를 바꾸지 않는다.
\begin{align}
\sigma_x\ket{u\,i}&=\ket{d\,i},
&
\sigma_x\ket{d\,i}&=\ket{u\,i},\\
\sigma_z\ket{u\,i}&=\ket{u\,i},
&
\sigma_z\ket{d\,i}&=-\ket{d\,i}.
\end{align}
두 번째 스핀에 대해서는 첫 번째 표지가 바뀌지 않는다.
\begin{align}
\tau_x\ket{a\,u}&=\ket{a\,d},
&
\tau_x\ket{a\,d}&=\ket{a\,u},\\
\tau_z\ket{a\,u}&=\ket{a\,u},
&
\tau_z\ket{a\,d}&=-\ket{a\,d}.
\label{eq:l6-spectator-action}
\end{align}
\(y\) 성분도 작용하는 자리를 뒤집으며,
\(\sigma_y\ket{u}=i\ket{d}\)와
\(\sigma_y\ket{d}=-i\ket{u}\)가 성립한다.

\begin{figure}[htbp]
\centering
\includegraphics[width=0.86\textwidth]{lecture_06_figure_05.png}
\caption{두 번째 스핀의 작용표. 첫 번째 지표는 바뀌지 않는다.}
\lectureframe{00:47:40}
\end{figure}

\begin{classroomqa}
\audiencequestion{\(\sigma_x\)를 적용하면 \(x\) 측정이 이루어지나요? 결과는 왜 중첩 상태가 아닌가요?}
\lecturerresponse{연산자의 작용 자체가 측정은 아닙니다. 행렬 \(\sigma_x\)를 \(z\) 기저벡터 \(\ket{u}\)에 곱하면 \(\ket{d}\)를 얻습니다. 이것은 고유값 방정식이 아닙니다. 고유벡터라면 연산자의 작용 뒤에도 방향이 유지되어야 합니다. 실제로 \(\ket{u}\)는 직교 상태로 변환되며, 이는 \(z\) 고유상태에 대한 \(x\) 측정이 최대한 무작위적인 이유를 보여 주는 또 다른 방식입니다. 고유벡터가 아닌 상태에 연산자를 작용시키는 것은 대개 중간 단계의 대수 계산이며, 그 물리적 의미는 계산하는 전체 물리량에 따라 정해집니다.}
\end{classroomqa}

서로 다른 인자에 작용하는 연산자들은 교환한다. 예를 들어,
\begin{equation}
\sigma_x\tau_x\ket{uu}
=\ket{dd}
=\tau_x\sigma_x\ket{uu}.
\end{equation}
같은 등식은 모든 복합계 기저벡터에서 성립하므로 모든 상태에서도 성립한다.
\begin{equation}
[\sigma_m,\tau_n]=0
\qquad\text{for all }m,n.
\label{eq:l6-local-commutation}
\end{equation}

\begin{figure}[htbp]
\centering
\includegraphics[width=0.86\textwidth]{lecture_06_figure_06.png}
\caption{서로 다른 자리에 작용하므로 교환자가 사라진다: \([\sigma_x,\tau_x]\).}
\lectureframe{00:56:40}
\end{figure}

\begin{classroomqa}
\audiencequestion{기저벡터 하나만 확인해서는 왜 두 연산자가 교환한다는 증명이 되지 않나요?}
\lecturerresponse{연산자가 같으려면 모든 벡터에 동일하게 작용해야 합니다. 기저의 모든 벡터에서 확인하면 충분한데, 선형성에 의해 그 결과가 모든 중첩 상태로 확장되기 때문입니다. 국소 연산자의 경우 네 기저 상태 모두에서 논증은 즉시 성립합니다. 각 연산자가 서로 다른 자리를 바꾸므로 작용 순서는 중요하지 않습니다.}
\end{classroomqa}

\begin{classroomqa}
\audiencequestion{두 스핀계 자체에 속하는 관측량도 있나요?}
\lecturerresponse{있습니다. \(\sigma_x\tau_x\), \(\sigma_x\tau_y\) 같은 곱과 그러한 곱들의 선형결합은 두 스핀의 공동 성질을 조사합니다. 국소 관측량은 각 인자의 성질을 기술하고, 상관 관측량은 두 인자가 어떻게 연관되어 있는지를 드러냅니다.}
\end{classroomqa}

\begin{classroomqa}
\audiencequestion{독립적으로 준비한 곱상태도 이미 얽혀 있나요?}
\lecturerresponse{아닙니다. 곱상태에서는 각 인자가 자체 상태벡터를 가지며, 국소 기댓값들은 인자별로 분해됩니다. 복합계 상태를 \(\sigma\)에 해당하는 상태 하나와 \(\tau\)에 해당하는 상태 하나로 분리할 수 없을 때 얽힘이 생깁니다.}
\end{classroomqa}

\section{싱글렛 상태}

최대로 얽힌 상태의 한 예는 다음과 같다.
\begin{equation}
\ket{S}
=\frac{1}{\sqrt2}\bigl(\ket{ud}-\ket{du}\bigr).
\label{eq:l6-singlet}
\end{equation}
두 항만 보아도 \(z\) 방향에서 완벽한 반상관이 드러난다. 한 스핀이 위쪽이면 다른 스핀은 아래쪽이다. 더 중요한 점은 어느 스핀에도 개별적인 순수 상태를 부여할 수 없다는 것이다.

\begin{figure}[htbp]
\centering
\includegraphics[width=0.86\textwidth]{lecture_06_figure_07.png}
\caption{이후 싱글렛으로 확인할 반대칭 두 스핀 상태.}
\lectureframe{01:09:25}
\end{figure}

한 스핀의 순수 상태에는 다음 식을 만족하는 단위벡터 \(\bm n\)가 항상 존재한다.
\begin{equation}
(\bm\sigma\cdot\bm n)\ket{\psi}=\ket{\psi}.
\end{equation}
다시 말해, 한 스핀의 순수 상태에서 블로흐 벡터 \begin{equation}
\bm r
=\bigl(\expect{\sigma_x},\expect{\sigma_y},\expect{\sigma_z}\bigr)
\end{equation}의 길이는 단위 길이이다. 따라서 세 성분이 모두 사라질 수는 없다. 상태가 주어지면 검출기를 \(\bm r\) 방향으로 정렬해 확실하게 \(+1\)를 얻을 수 있다. 익숙한 예로는 상태 \(\ket{u}\)에 대한 \(z\) 방향, 상태 \((\ket{u}+\ket{d})/\sqrt2\)에 대한 \(x\) 방향, 상태 \((\ket{u}+i\ket{d})/\sqrt2\)에 대한 \(y\) 방향 중 하나가 있다.

이제 싱글렛의 국소 평균을 계산하자.
\begin{align}
\expect{\sigma_z}_S
&=\frac12
(\bra{ud}-\bra{du})\sigma_z
(\ket{ud}-\ket{du})\\
&=\frac12
(\bra{ud}-\bra{du})
(\ket{ud}+\ket{du})\\
&=0.
\end{align}
\(\sigma_x\)에 대해서는,
\begin{equation}
\sigma_x\ket{S}
=\frac{\ket{dd}-\ket{uu}}{\sqrt2},
\end{equation}
이는 \(\ket{S}\)와 직교하므로 \(\expect{\sigma_x}_S=0\)이다. \(\sigma_y\)가 더하는 위상은 직교성에 영향을 주지 않으므로,
\begin{equation}
\expect{\sigma_x}_S
=\expect{\sigma_y}_S
=\expect{\sigma_z}_S=0.
\label{eq:l6-zero-local-sigma}
\end{equation}
대칭성에 따라,
\begin{equation}
\expect{\tau_x}_S
=\expect{\tau_y}_S
=\expect{\tau_z}_S=0.
\label{eq:l6-zero-local-tau}
\end{equation}

\begin{figure}[htbp]
\centering
\includegraphics[width=0.86\textwidth]{lecture_06_figure_08.png}
\caption{칠판에는 사라지는 국소 스핀 평균과 계산을 마무리할 때 사용한 \(y\) 성분의 뒤집힘이 기록되어 있다.}
\lectureframe{01:14:20}
\end{figure}

두 스핀 전체의 상태는 순수하며 완전히 알려져 있지만, 각 부분은 국소적으로 편극되지 않았다. 순수 한 스핀 상태의 블로흐 벡터는 0일 수 없으므로 \(\ket{S}\)는 곱상태일 수 없다.

\begin{classroomqa}
\audiencequestion{보통의 곱 상태에서도 국소 편극이 0이 아닌 방향이 있을까?}
\lecturerresponse{그렇다. 한 스핀에 대한 증명은 해당 인자에 적용되고, 다른 인자는 부수적인 역할만 한다. 따라서 순수 곱 상태의 각 인자에는 단위 블로흐 벡터가 있다. 싱글렛은 두 스핀 모두에서 이러한 곱 상태의 성질을 만족하지 않는다.}
\end{classroomqa}

\begin{classroomqa}
\audiencequestion{두 스핀을 싱글렛 상태로 어떻게 준비하나요?}
\lecturerresponse{반평행 정렬을 선호하는 결합은 싱글렛을 최저 에너지 상태로 만들 수 있습니다. 두 스핀을 상호작용하도록 가까이 둔 다음 주변으로 에너지를 잃게 하면 싱글렛 상태로 이완될 수 있습니다. 이 설명은 상호작용에 따라 달라집니다. 등방성 반강자성 해밀토니안 \(H_{\rm int}=J\,\bm\sigma\cdot\bm\tau\)에서 \(J>0\)이면, 싱글렛은 세 트립렛 상태보다 낮은 에너지를 갖습니다.}
\end{classroomqa}

\begin{classroomqa}
\audiencequestion{그 이완에는 시간이 얼마나 걸리나요?}
\lecturerresponse{시간 척도는 스핀 사이의 거리, 상호 결합의 세기, 자기 모멘트, 그리고 에너지를 가져가는 전자기 환경과의 결합에 따라 달라집니다. 일반적으로 결합이 강할수록 이완은 빨라집니다.}
\end{classroomqa}

\begin{classroomqa}
\audiencequestion{반대칭 조합의 에너지가 대칭 조합보다 낮을 수 있는 이유는 무엇인가요?}
\lecturerresponse{스핀 \(\tfrac12\) 입자 두 개의 경우, 반대칭 상태의 총스핀은 0이고 대칭 상태 세 개는 총스핀 1인 트립렛을 이룹니다. \(+\bm\sigma\cdot\bm\tau\)에 비례하는 상호작용은 싱글렛에 고유값 \(-3\)을, 트립렛에 \(+1\)를 부여하므로, 결합상수가 양수이면 싱글렛의 에너지가 더 낮습니다. 다른 해밀토니안에서는 준위의 순서가 다를 수 있습니다.}
\end{classroomqa}

\subsection{모든 스핀 기저에서 같은 상태}

더 편리한 국소 기저를 택한다고 해서 얽힘이 사라지지는 않는다. 뒤에서 다룰 기저 독립적인 척도로는 얽힘 엔트로피가 있다. 싱글렛은 그보다 더 강한 대칭성을 보인다. 다음을 사용하면,
\begin{equation}
\ket{u}=\frac{\ket{R}+\ket{L}}{\sqrt2},
\qquad
\ket{d}=\frac{\ket{R}-\ket{L}}{\sqrt2},
\end{equation}
이를 식~\eqref{eq:l6-singlet}에 대입하여 다음을 얻는다.
\begin{equation}
\ket{S}
=\frac{\ket{LR}-\ket{RL}}{\sqrt2}.
\label{eq:l6-singlet-x-basis}
\end{equation}
공통의 임의 정규직교 스핀 기저 \(\{\ket{\rm in},\ket{\rm out}\}\)에서도,
\begin{equation}
\ket{S}
=\frac{\ket{{\rm in}\,{\rm out}}-\ket{{\rm out}\,{\rm in}}}{\sqrt2},
\label{eq:l6-singlet-any-basis}
\end{equation}
전체 위상은 무관하므로 이를 제외한 결과다. 공통 회전 아래에서 세 트립렛 상태는 서로 섞이지만, 1차원 싱글렛 부분공간은 자기 자신으로만 변환된다.

\begin{figure}[htbp]
\centering
\includegraphics[width=0.86\textwidth]{lecture_06_figure_09.png}
\caption{위쪽-아래쪽 기저에서 오른쪽-왼쪽 기저 및 다른 공통 스핀 기저로 바꾸어도 반대칭 형태는 유지된다.}
\lectureframe{01:21:00}
\end{figure}

\begin{classroomqa}
\audiencequestion{얽힘은 상태를 쓰는 데 사용한 기저가 만들어 낸 착시일 뿐인가요?}
\lecturerresponse{아닙니다. 국소 기저를 바꾸면 계수는 달라지지만 상태가 인수분해되는지는 달라지지 않습니다. 얽힘 엔트로피는 이 불변 성질을 정량화합니다. 싱글렛에서는 공통 회전을 통해 이 불변성이 특히 뚜렷하게 드러납니다. 모든 스핀 기저에서 반대칭 형태가 그대로 재현됩니다.}
\end{classroomqa}

\section{상관 관측량}

국소 평균은 사라지므로 싱글렛의 정보는 공동 관측량에 담겨 있어야 한다. 다음에서 시작하자.
\begin{equation}
\sigma_z\tau_z\ket{ud}=-\ket{ud},
\qquad
\sigma_z\tau_z\ket{du}=-\ket{du}.
\end{equation}
따라서
\begin{equation}
\sigma_z\tau_z\ket{S}=-\ket{S},
\qquad
\expect{\sigma_z\tau_z}_S=-1.
\label{eq:l6-zz-correlation}
\end{equation}
이는 단순한 평균 결과가 아니다. 싱글렛은 이 곱 연산자의 고유상태이므로, 같은 축을 따라 \(z\)을 측정하면 부호는 항상 서로 반대다.

두 \(x\) 연산자는 각각 자신의 자리를 뒤집는다.
\begin{equation}
\sigma_x\tau_x\ket{ud}=\ket{du},
\qquad
\sigma_x\tau_x\ket{du}=\ket{ud}.
\end{equation}
두 항이 서로 바뀌면서 반대칭 상태의 부호가 뒤집힌다.
\begin{equation}
\sigma_x\tau_x\ket{S}=-\ket{S},
\qquad
\expect{\sigma_x\tau_x}_S=-1.
\end{equation}
\(y\) 연산자에서 생기는 두 위상도 같은 방식으로 상쇄되어 다음을 얻는다.
\begin{equation}
\expect{\sigma_y\tau_y}_S=-1.
\end{equation}
이와 달리 서로 다른 데카르트 성분들의 기댓값은 0이다. 예를 들면,
\begin{equation}
\sigma_x\tau_z\ket{S}
=-\frac{\ket{dd}+\ket{uu}}{\sqrt2},
\qquad
\expect{\sigma_x\tau_z}_S=0.
\end{equation}
이 모든 결과는 다음과 같이 요약된다.
\begin{equation}
\expect{\sigma_m\tau_n}_S=-\delta_{mn},
\qquad
\expect{(\bm\sigma\cdot\bm a)(\bm\tau\cdot\bm b)}_S
=-\bm a\cdot\bm b.
\label{eq:l6-singlet-correlation-tensor}
\end{equation}

\begin{figure}[htbp]
\centering
\includegraphics[width=0.86\textwidth]{lecture_06_figure_10.png}
\caption{상관 계산에서 얻은 \(\expect{\sigma_x\tau_x}=-1\)와 서로 다른 축 성분에 대한 결과 \(\expect{\sigma_x\tau_z}=0\).}
\lectureframe{01:32:20}
\end{figure}

\begin{classroomqa}
\audiencequestion{두 관측량의 곱은 물리적으로 무엇을 뜻하나요? 측정을 차례로 하나씩 하는 건가요?}
\lecturerresponse{일반적으로 에르미트 연산자의 곱은 복합계의 또 다른 관측량입니다. 여기서는 모든 \(\sigma_m\)가 모든 \(\tau_n\)와 교환하므로, 두 국소 성분을 동시에 측정해 그 수치 결과를 곱할 수 있습니다. 곱 연산자는 그 공동 결과를 나타내며, 측정 순서를 정할 필요는 없습니다.}
\end{classroomqa}

곱상태에서는 공동 기댓값이 인자별 기댓값의 곱으로 분해된다.
\begin{equation}
\expect{\sigma_m\tau_n}_{\rm prod}
=\expect{\sigma_m}_{\sigma}\expect{\tau_n}_{\tau}.
\end{equation}
싱글렛에서는 이 분해가 최대한 깨진다. 우변의 각 인자는 0이지만, 좌변의 같은 축 성분의 곱은 \(-1\)이다. 각 스핀은 개별적으로 완전히 무작위이지만, 공동으로 보면 같은 성분은 확실히 서로 반대다.

\begin{classroomqa}
\audiencequestion{두 스핀 사이의 모든 상호작용은 얽힘을 일으키나요?}
\lecturerresponse{일반적인 결합은 처음에 인수분해되어 있던 상태를 얽히게 할 수 있으며, 흔한 상호작용 중에도 그런 경우가 많습니다. 그렇다고 문자 그대로 반드시 그런 것은 아닙니다. 특별히 선택한 곱 고유상태는 위상만 얻을 수 있고, 국소 항만으로 이루어진 해밀토니안은 인자들을 얽히게 할 수 없습니다. 중요한 것은 곱상태로 보존되지 않는 상태에 비자명한 결합이 작용하는가입니다.}
\end{classroomqa}

\section{상관관계에서 벨의 질문으로}

완벽한 반상관 자체는 역설적이지 않다. 동전 하나와 10센트 동전을 주머니에 넣고, 보지 않은 채 하나씩 앨리스와 밥에게 나누어 준 다음 두 사람을 멀리 떨어뜨린다고 하자. 앨리스가 자신의 동전을 확인하면 밥의 동전도 즉시 알 수 있다. 상관관계는 동전들이 함께 있을 때 이미 정해졌으므로, 한쪽 결과를 알게 되어도 다른 위치로 신호가 전달되지는 않는다.

싱글렛에도 이런 단순한 특징이 있지만, 중요한 자유도가 더해진다. 앨리스와 밥은 측정 방향을 임의로 선택할 수 있고, 양자역학은 연속적으로 변하는 상관관계 \(-\bm a\cdot\bm b\)를 예측한다. 벨 정리는 이 모든 상관관계가 입자들이 미리 지니고 온 국소적인 고전적 지시로 설명될 수 있는지를 묻는다.

\begin{classroomqa}
\audiencequestion{벨 정리에서 사용하는 상태가 이 상태인가요?}
\lecturerresponse{그렇습니다. 벨 유형의 논증에서는 확정된 곱상태 \(\ket{ud}\)가 아니라 싱글렛 같은 얽힌 쌍을 사용합니다. 이 차이가 중요한 이유는 곱상태는 서로 분리된 국소 기술을 제공하는 반면, 싱글렛은 모든 공통 기저에 걸친 분리 불가능한 상관 구조를 담기 때문입니다.}
\end{classroomqa}

\begin{classroomqa}
\audiencequestion{같은 축 측정 결과는 항상 반대인데 왜 반상관을 통계적이라고 하나요?}
\lecturerresponse{각각의 국소 결과는 무작위입니다. 앨리스는 \(+1\)와 \(-1\) 중 어떤 결과를 얻을지 예측할 수 없습니다. 조건부 관계는 정확합니다. 같은 축 측정에서 한쪽 결과를 알면 다른 쪽 결과는 그 반대값으로 정해집니다. 상관관계는 같은 축 성분의 곱에 불확실성이 있다는 뜻이 아니라, 공동 확률분포를 가리킵니다.}
\end{classroomqa}

이 문제를 구체적으로 살펴보는 유용한 방법은 시뮬레이션이다. 고전 컴퓨터 한 대로 스핀 하나에 대한 모든 실험을 모사할 수 있다. 상태를 저장하고, Schr\"odinger 방정식에 따라 시간에 따라 변화시키고, 요청된 확률을 계산한 다음, 난수 생성기로 결과 하나를 선택하면 된다. 반복 측정과 자기장에 의한 세차도 모두 재현할 수 있다.

이제 얽힌 쌍을 고전 컴퓨터 두 대에 나누어 주자. 쌍이 준비되는 동안 두 컴퓨터가 상호작용하도록 한 뒤 연결을 끊고, 멀리 떨어진 각 위치에서 측정 방향을 독립적으로 선택한다. 미리 정해 둔 공유 난수를 사용하더라도, 분리된 국소 고전 프로그램 두 개로 모든 측정 방향 선택에 대한 양자 상관관계를 재현할 수는 없다. 전역적인 고전 시뮬레이션은 가능하지만, 공동 상태에 계속 접근할 수 있어야 하거나 다른 방식으로 두 측면을 조정해야 한다. 이는 국소 숨은 변수 모형이 아니다.

국소적인 고전 장부 맞추기가 실패한다고 해서 광속보다 빠른 통신 경로가 생기는 것은 아니다. 앨리스의 국소 측정 결과는 여전히 무작위이므로, 나중에 기록을 서로 대조하지 않고 밥이 읽을 수 있는 메시지를 앨리스가 제어할 수 없다. 비국소성에 관한 용어는 여러 방식으로 쓰이지만, 조작적 구분은 명확하다. 벨 상관관계는 국소 고전적 설명을 배제하면서 무신호성을 보존한다.

이로써 복합계 상태의 구성에서 새로운 것이 드러났다. 텐서곱은 국소적으로 이전의 한 스핀 물리를 보존하지만, 그 공간의 일반적인 벡터는 각 스핀의 상태로 나눌 수 없다. 싱글렛에서는 모든 국소 편극이 사라지고, 공동 스핀 성분은 완벽하고 기저에 무관한 반상관을 지닌다. 다음 단계에서는 이 구조를 바탕으로 고전적 추론을 시험하는 벨의 정량적 검사를 세운다.
